\songtitle{Entre mundos}

\tag*{2026}\tag{folk-rock}\tag{original-lyrics}\tag{spanish-lyrics}
\begin{notes}
Written by Carlos Thompson -- a lifelong habit of building alternate worlds as an escape from daily friction (traffic, bureaucracy, the gap between how things are and how they could be), and how generative AI became one more tool for exploring them. Retrieved from Wyomee Dank's Suno page (V6-MINI, 18 September 2026).
\end{notes}
\begin{style}
Rock with folk delivery, drum kit, ringing electric guitars and grounded bass, fast-paced driving groove, clean small-venue room sound with close-mic definition, mellow male vocals in a conversational folk cadence
\end{style}
\begin{lyrics}
\annotation{Estrofa 1}\\
Caminar la ciudad, conducir en tráfico de locos,\\
un sistema de transporte público entre tumultos o sofocado en el mismo trancón,\\
y saber que pudo haber sido diferente, administraciones con otros focos\\
o el escape a la fricción con mi propia imaginación.

\annotation{Estrofa 2}\\
Imaginarme cómo sería ese mundo ideal, imaginarme viviendo en él.\\
Planes elaborados desde lo grande en los sistemas hasta lo cotidiano.\\
Las diferencias en la fricción entre este mundo y el alterno en mi propia piel,\\
y pensar cómo pudo haber sido si pudiera hacerlo con mi propia mano.

\annotation{Puente}\\
La fantasía de poder sin dejar de ser quien soy.\\
No es por imaginarme ser el centro de atención.\\
Es más haber logrado ese otro mundo hoy\\
y vivir mi vida, pero así, sin fricción.

\annotation{Estribillo}\\
Entre mundos alternos y en mi propia vida\\
me quedo pensando más de lo prudente,\\
archivos y hojas que sobreviven la barrida,\\
lo que albergo aún en mi propia mente.

Entre mundos reales y los imaginarios\\
me quedo pensando sin poder evitarlo,\\
eventos mundanos y extraordinarios,\\
lo que sea si es que puedo pensarlo.

\annotation{Estrofa 3}\\
Leer un artículo en línea sobre cómo escribir y ahora quiero escribir.\\
Estudiar un nuevo lenguaje de programación y pensar cómo sería mi propio lenguaje.\\
Historia, biología, veo lo que hay y una parte de mí quiere vivir\\
en nuevas realidades, imaginar qué sería si, sea o no mi alterno viaje.

\annotation{Estrofa 4}\\
Llevo inventando mundos desde que me alcanza la memoria.\\
Jugar a ser quien no soy, pero también un yo diferente u otra ficción.\\
Otras lenguas, otros mundos, otros cuentos y también otra historia,\\
y buscando formas de poder materializar mi visión.

\annotation{Puente}\\
Llegan los modelos de lenguaje e IAs generativas,\\
nuevas formas de explorar el mundo,\\
asistentes y editores de mis narrativas.\\
Suficiente, si bien no tan profundo.

\annotation{Estribillo}\\
Entre mundos alternos y en mi propia vida\\
me quedo pensando más de lo prudente,\\
archivos y hojas que sobreviven la barrida,\\
lo que albergo aún en mi propia mente.

Entre mundos reales y los imaginarios\\
me quedo pensando sin poder evitarlo,\\
eventos mundanos y extraordinarios,\\
lo que sea si es que puedo pensarlo.

\annotation{Desglose}\\
Y sin darme cuenta tengo ya una novela\\
con todo: escenas, retratos, banda sonora.\\
He explorado mi mente, lo que teme y anhela,\\
y un universo alterno que no se demora.

He cantado a mis miedos, mis desvaríos,\\
a mis mundos alternos y mi fantasía,\\
a las mujeres reales que pensé mis amoríos,\\
a las mujeres imaginarias que dictó a la IA.

\annotation{Estribillo}\\
Entre mundos alternos y en mi propia vida\\
me quedo pensando más de lo prudente,\\
archivos y hojas que sobreviven la barrida,\\
lo que albergo aún en mi propia mente.

Entre mundos reales y los imaginarios\\
me quedo pensando sin poder evitarlo,\\
eventos mundanos y extraordinarios,\\
lo que sea si es que puedo pensarlo.

Entre mundos alternos,\\
entre mundos reales,\\
mi propia vida\\
y lo imaginario.
\end{lyrics}

